% GENERATED FILE. Edit canonical lessons, then run npm run generate:books.
% canonical-source-hash: fnv1a64:803ea1e48e83dd93
% canonical-lessons: AR-C287-kayfa-asilu, AR-C287-idhhab, AR-R287-asking-the-way

\chapter{Asking The Way}
\label{ch:ar-asking-the-way}
\hlchaptermodality{287}

\begin{chapteropening}
\textbf{By the end of this chapter:} \emph{I can ask how to get somewhere and follow directions given as orders.}

Complete AR-R287-asking-the-way: ask the way with kayfa aṣilu ilā and give the answer as orders.
\end{chapteropening}

\section[\texorpdfstring{\ar{كيف} \ar{أصل} \ar{إلى} (kayfa aṣilu ilā)}{كيف أصل إلى (kayfa aṣilu ilā)}]{\ar{كيف} \ar{أصل} \ar{إلى} (kayfa aṣilu ilā) — how do I get to …?}

\label{lesson:AR-C287-kayfa-asilu}



\begin{quote}
Say \textquotedblleft{}I will travel tomorrow\textquotedblright{}. (\emph{Sa-usāfiru ghadan}.) Then \textquotedblleft{}I want to go to the market\textquotedblright{}. (\emph{Urīdu an adhhaba ilā as-sūq}.)

You are in a city you do not know. This is what you say to the first person who stops.
\end{quote}

\subsection*{You'll want to know: \ar{كيف} \ar{أصل} \ar{إلى}}
\textbf{\ar{كيف} \ar{أصل} \ar{إلى}} — \emph{kayfa aṣilu ilā} — \textquotedblleft{}how do I get to\textquotedblright{}, word for word \textquotedblleft{}how do I arrive at\textquotedblright{}.

\begin{quote}
\textbf{\ar{من} \ar{فضلك،} \ar{كيف} \ar{أصل} \ar{إلى} \ar{المحطة؟}} — \emph{min faḍlik, kayfa aṣilu ilā al-maḥaṭṭa?} — Excuse me, how do I get to the station?
\end{quote}

\emph{Aṣilu}, I arrive, is the present of \emph{waṣala}, and the root's \emph{w} has dropped, as it did in \emph{yajibu}. To a woman, say \emph{min faḍliki}.

\begin{grammarlens}[title={Grammar Lens: the questions around it}]
Two shorter questions often come with it. \textbf{\ar{أين}} — \emph{ayna}, \textquotedblleft{}where\textquotedblright{}: \emph{ayna al-matḥaf?}, where is the museum? And \textbf{\ar{هل}} with \emph{baʿīd} or \emph{qarīb}: \emph{hal al-matḥaf baʿīd?}, is the museum far? The answer may be \emph{qarīb min hunā}, near here, or \emph{baʿīd qalīlan}, a little far.

The answer to \emph{kayfa aṣilu} is a set of orders, which the next lesson builds.
\end{grammarlens}

\begin{cousinweb}
\textbf{\ar{أصل}} (\emph{aṣilu}) is built on the native Arabic root \emph{w-ṣ-l}, \textquotedblleft{}to reach, to join\textquotedblright{} (Lane), the root of \emph{waṣala}, \textquotedblleft{}he arrived\textquotedblright{}, and of \emph{ittaṣala}, \textquotedblleft{}he phoned\textquotedblright{}, a call that joins two people. \textbf{\ar{كيف}} (\emph{kayfa}) and \textbf{\ar{إلى}} (\emph{ilā}) are inherited particles, not built on three-letter roots.
\end{cousinweb}

\subsection*{Your turn}
Say these aloud:

\begin{itemize}
\raggedright
  \item \emph{min faḍlik, kayfa aṣilu ilā al-maḥaṭṭa?}
  \item the same, to a woman, for the museum — \emph{min faḍliki, kayfa aṣilu ilā al-matḥaf?}
  \item \textquotedblleft{}Is the museum far?\textquotedblright{} — \emph{hal al-matḥaf baʿīd?}
\end{itemize}

\begin{tcolorbox}[breakable,colback=teal!4,colframe=teal!35!black,title={Before you move on}]
\textquotedblleft{}How do I get to the station?\textquotedblright{} (\emph{\textbf{Kayfa aṣilu ilā al-maḥaṭṭa?}}) Which verb is inside \emph{aṣilu}? (\emph{\textbf{Waṣala}}, to arrive.) Which letter of its root has dropped? (\textbf{The \emph{w}.})
\end{tcolorbox}
\section[\texorpdfstring{\ar{اذهب} \ar{مستقيما} (idhhab mustaqīman)}{اذهب مستقيما (idhhab mustaqīman)}]{\ar{اذهب} \ar{مستقيما} (idhhab mustaqīman) — go straight on}

\label{lesson:AR-C287-idhhab}



\begin{quote}
Ask the way to the station. (\emph{Min faḍlik, kayfa aṣilu ilā al-maḥaṭṭa?}) Then ask someone to say it again. (\emph{Karrir min faḍlik}.)

\emph{Karrir} was an order. The answer to your question will be a string of them.
\end{quote}

\begin{grammarlens}[title={Grammar Lens: an order from the present}]
Start from \textquotedblleft{}you do it\textquotedblright{}: \emph{tadhhabu}, you go. Take off the \emph{ta-} and the final vowel, and \emph{-dhhab} is left, which cannot open with two consonants. Arabic puts a short vowel in front: \emph{u-} when the stem's vowel is \emph{u}, and \emph{i-} otherwise.

\begin{quote}
\emph{tadhhabu} $\to$ \textbf{\ar{اذهب}} \emph{idhhab}, go · \emph{taʿburu} $\to$ \textbf{\ar{اعبر}} \emph{uʿbur}, cross · \emph{taktubu} $\to$ \textbf{\ar{اكتب}} \emph{uktub}, write
\end{quote}

To a woman, add \emph{-ī}: \emph{idhhabī}, \emph{uʿburī}. One common order is short and irregular: \textbf{\ar{خذ}} \emph{khudh}, take, from \emph{akhadha}; to a woman, \emph{khudhī}.
\end{grammarlens}

\begin{grammarlens}[title={Grammar Lens: the way, in orders}]
\begin{quote}
\textbf{\ar{اذهب} \ar{مستقيما،} \ar{ثم} \ar{خذ} \ar{الشارع} \ar{الأول} \ar{على} \ar{اليمين}.} — \emph{idhhab mustaqīman, thumma khudh ash-shāriʿa al-awwala ʿalā al-yamīn} — Go straight on, then take the first street on the right.
\end{quote}

\begin{quote}
\textbf{\ar{اعبر} \ar{الجسر،} \ar{والمحطة} \ar{على} \ar{اليسار}.} — \emph{uʿbur al-jisr, wa-l-maḥaṭṭa ʿalā al-yasār} — Cross the bridge, and the station is on the left.
\end{quote}

\emph{Mustaqīman}, with \emph{-an}, says how you go. \emph{ʿAlā al-yamīn} and \emph{ʿalā al-yasār} say where a thing is. \emph{Thumma} strings the steps together, as it did in a story.
\end{grammarlens}

\subsection*{Your turn}
Say these aloud:

\begin{itemize}
\raggedright
  \item \emph{tadhhabu} … \emph{idhhab}; \emph{taʿburu} … \emph{uʿbur}
  \item \textquotedblleft{}Go straight on\textquotedblright{} — \emph{idhhab mustaqīman}
  \item \textquotedblleft{}Take the first street on the right\textquotedblright{} — \emph{khudh ash-shāriʿa al-awwala ʿalā al-yamīn}
  \item \textquotedblleft{}Cross the bridge\textquotedblright{}, to a woman — \emph{uʿburī al-jisr}
  \item the whole answer — \emph{idhhab mustaqīman, thumma uʿbur al-jisr, wa-l-maḥaṭṭa ʿalā al-yasār}
\end{itemize}

\begin{tcolorbox}[breakable,colback=teal!4,colframe=teal!35!black,title={Before you move on}]
Where does an order come from? (\textbf{The present}, with \emph{ta-} taken off.) \textquotedblleft{}Go\textquotedblright{}? (\emph{\textbf{Idhhab.}}) \textquotedblleft{}Take\textquotedblright{}? (\emph{\textbf{Khudh}}, the short one.) \textquotedblleft{}On the left\textquotedblright{}? (\emph{\textbf{ʿAlā al-yasār.}}) What do you add for a woman? (\emph{\textbf{-ī}}: \emph{idhhabī}.)
\end{tcolorbox}
\section[(review)]{Review — asking the way, and following it}

\label{lesson:AR-R287-asking-the-way}



\begin{quote}
Say \textquotedblleft{}I must go now\textquotedblright{} and \textquotedblleft{}you ought to rest\textquotedblright{}. (\emph{Yajibu an adhhaba al-ān}, \emph{yanbaghī an tastarīḥa}.) Then \textquotedblleft{}I won't go, because of the rain\textquotedblright{}. (\emph{Lan adhhaba, bi-sabab al-maṭar}.)
\end{quote}

\begin{grammarlens}[title={What you've built}]
Two lessons, and both halves of a conversation in the street. Your half: \emph{min faḍlik, kayfa aṣilu ilā …?}, with \emph{ayna} and \emph{hal … baʿīd?} beside it. Their half: orders, the present with its \emph{ta-} taken off, strung together with \emph{thumma}.

One exchange, whole:

\begin{quote}
— \emph{min faḍlik, kayfa aṣilu ilā al-funduq?} — \emph{idhhab mustaqīman, thumma khudh ash-shāriʿa al-awwala ʿalā al-yasār.} — \emph{hal huwa baʿīd?} — \emph{lā, qarīb min hunā.}
\end{quote}
\end{grammarlens}

\subsection*{Your turn}
Say these aloud:

\begin{itemize}
\raggedright
  \item \textquotedblleft{}Excuse me, how do I get to the hotel?\textquotedblright{} — \emph{min faḍlik, kayfa aṣilu ilā al-funduq?}
  \item \textquotedblleft{}Where is the museum?\textquotedblright{} — \emph{ayna al-matḥaf?}
  \item \textquotedblleft{}Go straight on, then take the first street on the left\textquotedblright{}
  \item \textquotedblleft{}The station is at the corner\textquotedblright{} — \emph{al-maḥaṭṭa ʿinda az-zāwiya}
  \item the way from your door to the centre of your own city, in three orders
\end{itemize}

\begin{tcolorbox}[breakable,colback=teal!4,colframe=teal!35!black,title={Before you move on}]
\textquotedblleft{}How do I get to …?\textquotedblright{} (\emph{\textbf{Kayfa aṣilu ilā …?}}) \textquotedblleft{}Cross\textquotedblright{}? (\emph{\textbf{Uʿbur.}}) \textquotedblleft{}Take\textquotedblright{}, to a woman? (\emph{\textbf{Khudhī.}}) \textquotedblleft{}On the right\textquotedblright{}? (\emph{\textbf{ʿAlā al-yamīn.}})
\end{tcolorbox}
